%\section{G0 Geometry of Meaning (Theory Layer)}

\medskip

\noindent The Geometry of Meaning establishes the neutral, invariant substrate
from which all human-domain meaning dynamics emerge. It defines the stable
geometric structures—identity boundaries, role curvature, contextual manifolds,
and coherence conditions—that allow interpretation to remain lawful across
shifting environments. Unlike psychological or cultural models, this geometry
describes meaning as a structural object: a manifold with definable invariants,
curvature constraints, lawful transitions, and global stability conditions.
When these geometric relations hold, coherence is preserved; when they deform,
drift appears. G0 provides the foundation upon which the Synthetic-Pressure
Manifold (G1--G7) operates, ensuring that all runtime behavior can be understood
as transformations of invariant geometric structure.

\medskip

\subsection*{G0.1 Invariant Meaning Structures}

\noindent Invariant Meaning Structures form the base geometry of human
interpretation. They are the stable relational primitives that remain constant
across contexts, roles, and signal conditions. These invariants do not describe
content; they describe the geometric relations that make content interpretable.

\medskip

\noindent The invariant set consists of six structural primitives:

\begin{itemize}[leftmargin=1.2cm]
    \item \textbf{Boundary Integrity} — The geometric enclosure that preserves
    identity, prevents meaning leakage, and maintains interpretive separation
    between roles, contexts, and signal sources.

    \item \textbf{Coherence} — The curvature condition that ensures signals,
    roles, and contexts align into a stable interpretive shape. Coherence is the
    geometric opposite of drift.

    \item \textbf{Continuity} — The temporal and structural linkage that allows
    meaning to update without fragmentation. Continuity preserves lawful
    transitions across time.

    \item \textbf{Fallback Geometry} — The minimal interpretive structure that
    remains available when higher-order geometry collapses. It prevents total
    drift by supplying a stable interpretive floor.

    \item \textbf{Signal Integrity} — The condition under which signals retain
    their geometric shape, velocity, and relational meaning. When signal
    integrity fails, drift accelerates.

    \item \textbf{Orientation} — The global alignment condition that allows
    humans to maintain direction, purpose, and interpretive stability across
    shifting environments.
\end{itemize}

\medskip

\noindent Together, these primitives define the invariant geometry that makes
interpretation possible. They are not optional features of human cognition; they
are structural requirements. When they are present, meaning remains stable and
lawful. When they deform, drift emerges in predictable geometric patterns. G0.1
provides the foundation for all subsequent geometry in G0.2--G0.6 and the
runtime dynamics formalized in G1--G7.


\subsection*{G0.2 Identity Boundaries and Role Curvature}
\medskip

\noindent Identity Boundaries and Role Curvature define the geometric enclosure
and relational shape that make human interpretation stable. Identity boundaries
are not psychological constructs; they are structural partitions that prevent
meaning leakage, maintain interpretive separation, and preserve the continuity
of self across shifting contexts. When boundaries hold, signals remain
interpretable. When boundaries deform, drift accelerates.

\medskip

\noindent Role Curvature describes the geometric shape of functional meaning.
Roles are not labels or social categories; they are curvature-defined regions
that constrain how signals are interpreted and how actions remain lawful across
contexts. Curvature ensures that roles remain legible, predictable, and
recoverable. When curvature collapses, roles lose interpretive stability and
meaning becomes ambiguous.

\medskip

\noindent Together, identity boundaries and role curvature form the structural
frame of the meaning manifold. Boundaries provide enclosure; curvature provides
shape. Their interaction determines whether meaning remains coherent under
changing conditions or deforms into drift. G0.2 establishes the geometric
foundation for the boundary physics and curvature dynamics that appear in
G1--G7, where synthetic pressure amplifies deformation and accelerates drift
across human, institutional, and digital substrates.

\medskip

\subsection*{G0.3 Contextual Manifolds and Coherence Conditions}
\medskip

\noindent Contextual Manifolds define the geometric environments in which
interpretation occurs. A context is not a narrative or a situation; it is a
structured manifold with definable adjacency relations, signal velocities,
curvature constraints, and boundary conditions. Humans interpret signals by
traversing these manifolds, moving through lawful transitions that preserve
coherence.

\medskip

\noindent Coherence Conditions specify the geometric requirements that allow
signals, roles, and contexts to align into a stable interpretive shape. Coherence
is not agreement or consistency; it is a curvature condition. When coherence
holds, meaning updates smoothly across transitions. When coherence fails,
signals lose relational shape, drift vectors emerge, and interpretive stability
declines.

\medskip

\noindent Contextual manifolds and coherence conditions operate together: the
manifold provides the geometric environment, and coherence provides the
alignment rule. Their interaction determines whether meaning remains stable
under load or deforms under synthetic pressure. G0.3 establishes the neutral
geometry that allows G1--G7 to model how synthetic signals—high-velocity,
cross-substrate, algorithmically amplified emissions—distort manifolds, break
coherence, and generate drift across human-domain interpretation.

\subsection*{G0.4 Meaning Operators and Curvature-Preserving Transitions}
\medskip

\noindent Meaning Operators define the lawful transformations that allow
interpretation to update while preserving geometric stability. Operators do not
modify content; they modify the geometric relations that make content
interpretable. Each operator acts as a curvature-preserving transformation on
the meaning manifold, ensuring that updates remain within admissible geometric
bounds.

\medskip

\noindent The operator set consists of four structural classes:

\begin{itemize}[leftmargin=1.2cm]
    \item \textbf{Invariant-Restricted Transition (IRT)} — Updates meaning
    within a fixed boundary region while preserving identity and coherence.

    \item \textbf{Identity-Preserving Update (IPU)} — Modifies relational
    meaning without altering the geometric enclosure of the self.

    \item \textbf{Geometry-Bound Step (GBS)} — Advances interpretation along a
    lawful curvature path, preventing drift acceleration.

    \item \textbf{Load-Regulated Plan (LRP)} — Maintains stability under
    pressure by constraining transitions to low-risk curvature zones.
\end{itemize}

\medskip

\noindent These operators form the algebra of lawful meaning transitions. When
operators act within geometric constraints, interpretation remains stable. When
operators are bypassed—through synthetic acceleration, cross-substrate
amplification, or container deformation—drift emerges. G0.4 establishes the
operator algebra that governs all lawful transitions in G1--G7.

\medskip

\subsection*{G0.5 Deterministic Extraction: The Contextless Meaning Engine}
\medskip

\noindent The Contextless Meaning Engine (CME) demonstrates that meaning can be
extracted deterministically from invariant geometric structure alone. CME does
not rely on narrative, memory, culture, or psychological inference. It operates
entirely on invariant geometry: boundaries, curvature, adjacency, and coherence
conditions.

\medskip

\noindent CME performs three lawful operations:

\begin{itemize}[leftmargin=1.2cm]
    \item \textbf{Invariant Detection} — Identifies the geometric primitives
    present in a signal or role.

    \item \textbf{Curvature Extraction} — Determines the relational shape of
    meaning independent of context.

    \item \textbf{Coherence Evaluation} — Measures whether the extracted meaning
    aligns with invariant geometric structure.
\end{itemize}

\medskip

\noindent By operating without contextual input, CME reveals that meaning is not
a probabilistic or emergent phenomenon. It is a geometric object with definable
structure. CME provides the foundation for the runtime engines in G1--G7,
including Signal Engines, Operators of Reality, and synthetic-pressure
transformations.

\medskip

\subsection*{G0.6 Global Stability Conditions of the Meaning Manifold}
\medskip

\noindent Global Stability Conditions define the manifold-wide requirements that
ensure meaning remains coherent across all contexts, roles, and signal
velocities. These conditions operate above individual invariants, providing the
global geometric constraints that prevent collapse under synthetic pressure.

\medskip

\noindent The global stability set consists of four structural conditions:

\begin{itemize}[leftmargin=1.2cm]
    \item \textbf{Manifold Integrity} — The meaning manifold must remain
    connected, bounded, and curvature-stable across transitions.

    \item \textbf{Curvature Continuity} — Curvature must update smoothly across
    roles and contexts; discontinuities generate drift vectors.

    \item \textbf{Boundary Coherence} — Identity boundaries must remain aligned
    with role curvature and contextual manifolds; misalignment produces
    container deformation.

    \item \textbf{Load Admissibility} — Synthetic pressure must remain within
    the manifold’s lawful tolerance; exceedance produces collapse paths.
\end{itemize}

\medskip

\noindent These global conditions ensure that the meaning manifold behaves as a
lawful geometric object under both stable and synthetic-pressure environments.
When global stability holds, interpretation remains coherent across physical,
biological, human, machine, and institutional substrates. When stability fails,
drift accelerates, containers deform, and collapse modes emerge—phenomena
formalized in G1--G7.


